EEseshq1BnC
The selected region contains parts (B) and (C) of Question One, which cover theoretical and practical aspects of network theorems, inductors, and capacitors.
Here are the complete, structured solutions for both sections:
Part (B): Theorem Comparisons
i. State the principle of each theorem
Thévenin's Theorem: Any linear, bilateral, two-terminal DC network can be replaced by an equivalent circuit consisting of a single voltage source ($V_{th}$) in series with a single resistor ($R_{th}$).
Norton's Theorem: Any linear, bilateral, two-terminal DC network can be replaced by an equivalent circuit consisting of a single constant current source ($I_n$) in parallel with a single resistor ($R_n$).
Superposition Theorem: In a linear, bilateral network containing multiple independent sources, the total response (current or voltage) in any branch is the algebraic sum of the individual responses caused by each independent source acting alone while all other independent sources are replaced by their internal resistances.
ii. Identify the type of circuits/conditions where each is most useful
Thévenin's Theorem: Most useful when analyzing circuits where a single load component ($R_L$) changes frequently while the rest of the network remains constant.
Norton's Theorem: Most useful for parallel network configurations and for designing or analyzing current-driven systems or electronic devices like transistors.
Superposition Theorem: Most useful when a circuit contains multiple power sources of different types (e.g., a mix of AC and DC sources, or sources operating at different frequencies).
iii. Explain how you would decide which theorem to apply in a given problem
Choose Thévenin’s if the goal is to find the voltage across or maximum power delivered to a variable series load.
Choose Norton’s if the goal is to determine the short-circuit current performance or analyze parallel-connected branches.
Choose Superposition when multiple independent sources are present and simpler nodal or mesh analyses are complicated, particularly if the sources cannot be combined directly due to differing characteristics or frequencies.
